\documentclass[11pt]{article}
\usepackage[a4paper,margin=2.6cm]{geometry}
\usepackage{amsmath,amssymb,amsthm}
\usepackage{booktabs}
\usepackage[colorlinks=true,linkcolor=blue,citecolor=blue,urlcolor=blue]{hyperref}

\newtheorem{theorem}{Theorem}
\newtheorem{proposition}[theorem]{Proposition}
\newtheorem{lemma}[theorem]{Lemma}
\theoremstyle{definition}
\newtheorem{heuristic}[theorem]{Heuristic}
\newtheorem{question}[theorem]{Question}
\theoremstyle{remark}
\newtheorem{remark}[theorem]{Remark}

\newcommand{\mz}{m_0}
\newcommand{\ord}{\operatorname{ord}}

\title{A counterexample to a conjecture of Zhi-Wei Sun\\ on sums of a prime and two powers of two}
\author{Yotam Moshe \and Idan Gan Levi \and Ohav Betito \and Boaz Cohen \and Shay Gez Akiva}
\date{25 November 2026 \quad (draft)}

\begin{document}
\maketitle

\begin{abstract}
In 2009 Zhi-Wei Sun conjectured that every odd integer greater than $34$ can be written as
$p+2^x+7\cdot 2^y$ with $p$ a prime congruent to $1$ modulo $6$ and $x,y$ positive integers
(OEIS A157218). We show that the conjecture is false: $445466268731$ has no such representation,
and it is the only counterexample below $10^{12}$. The proof of non-representability is a table
of $342$ explicit prime factors. For the sister conjecture with $p\equiv 5 \pmod 6$ (A157225),
disproved in 2026, we determine all four counterexamples below $10^{12}$. We describe the search
algorithm, a bit-parallel sieve that tests about $4\cdot 10^{9}$ integers per second on a desktop
computer, explain which arithmetic features make $445466268731$ exceptional, and give a heuristic
model that predicts the observed number of counterexamples and suggests that there are infinitely many.
\end{abstract}

\section{Introduction}

Romanoff proved in 1934 that a positive proportion of the integers are of the form $p+2^k$, and
Erd\H{o}s~\cite{erdos} used covering congruences to exhibit an arithmetic progression of odd numbers
none of which has this form. Allowing two powers of two, Crocker~\cite{crocker} showed that there are
still infinitely many odd integers not of the form $p+2^a+2^b$, and Pan~\cite{pan} gave a
quantitative version. In 2009 Sun~\cite{sun} asked what happens when one of the powers is
multiplied by a small odd constant, and recorded a family of conjectures in the OEIS~\cite{oeis},
among them the following two.

\medskip
\noindent\textbf{Conjecture A157218} (Sun, 24 February 2009). \emph{Every odd integer $m>34$ can be
written as $m=p+2^x+7\cdot 2^y$ with $p$ a prime, $p\equiv 1 \pmod 6$, and integers $x,y\ge 1$.}

\medskip
\noindent\textbf{Conjecture A157225} (Sun, 24 February 2009). \emph{Every odd integer $m>20$ other
than $25$, $31$ and $1983$ can be written as $m=p+2^x+7\cdot 2^y$ with $p$ a prime,
$p\equiv 5 \pmod 6$, and integers $x,y\ge 1$.}

\medskip
Both were verified for $m<1.5\cdot 10^{8}$ by Sun and Q.-H.~Hou. In 2026 Adamczewski~\cite{oeisopen}
reported, as part of a benchmark in which language models attempt open OEIS conjectures, that
A157225 fails at $m=1433987797$. Conjecture A157218 was not part of that benchmark and remained open.

\begin{theorem}\label{thm:main}
The odd integer
\[
  \mz \;=\; 445466268731 \;=\; 11^{3}\cdot 2999\cdot 111599
\]
cannot be written as $p+2^x+7\cdot 2^y$ with $p$ a prime congruent to $1$ modulo $6$ and $x,y\ge 1$.
Hence Conjecture A157218 is false.
\end{theorem}

\begin{theorem}[computational]\label{thm:search}
\begin{enumerate}
\item[(i)] $\mz$ is the only odd integer $m$ with $34<m<10^{12}$ that is not of the form
$p+2^x+7\cdot 2^y$ with $p\equiv 1 \pmod 6$ prime and $x,y\ge 1$.
\item[(ii)] The odd integers $m$ with $20<m<10^{12}$ that are not of the form $p+2^x+7\cdot 2^y$ with
$p\equiv 5\pmod 6$ prime and $x,y\ge 1$ are exactly
\[ 25,\quad 31,\quad 1983,\quad 1433987797,\quad 45483661477,\quad 85305551767,\quad 367084269651 . \]
\item[(iii)] The odd integers $m$ with $30<m<10^{12}$ that are not of the form $p+2^x+11\cdot 2^y$ with
$p\equiv 1\pmod 6$ prime and $x,y\ge 1$ are exactly $35$, $41$, $47$, $101$, $167$, $2021$ and $119171$,
in agreement with Sun's conjecture for OEIS A157237 (previously checked to $1.5\cdot 10^{8}$).
\end{enumerate}
\end{theorem}

Theorem~\ref{thm:main} has a short proof that can be checked by hand-sized arithmetic
(Section~\ref{sec:proof}). Theorem~\ref{thm:search} rests on a computer search; the algorithm and
the measures taken to validate it are described in Section~\ref{sec:algo}. In
Section~\ref{sec:structure} we ask why $\mz$ fails and how many counterexamples one should expect.
All programs and data are available at
\begin{center}\url{https://github.com/yotam1001/sun-a157218-counterexample}\end{center}

\section{Proof of Theorem~\ref{thm:main}}\label{sec:proof}

Suppose $\mz=p+2^x+7\cdot 2^y$ with $x,y\ge 1$ and $p\equiv 1\pmod 6$. Since $\mz\equiv 5\pmod 6$ we
get $2^x+7\cdot 2^y\equiv 4\pmod 6$, hence $2^x+2^y\equiv 1\pmod 3$. As $2^k\equiv 2$ or $1 \pmod 3$
according as $k$ is odd or even, this forces $x$ and $y$ to be both odd. Writing $x=2a+1$, $y=2b+1$,
\[
  p \;=\; \mz-2\cdot 4^{a}-14\cdot 4^{b},\qquad a,b\ge 0 .
\]
There are exactly $342$ pairs of odd $(x,y)$ with $2^x+7\cdot 2^y<\mz$ (namely $1\le x\le 37$,
$1\le y\le 35$ odd; all $19\cdot 18$ such pairs satisfy the inequality). For each of them
Table~\ref{tab:cert} in the appendix gives a prime $q<p$ dividing $p$, so $p$ is composite. \qed

\medskip
The least prime factors are small for most candidates: $5$ divides $90$ of the $342$ numbers, $7$
divides $81$ of the rest, $11$ a further $34$, and $13$ a further $13$; only $41$ candidates have
no prime factor below $500$. The largest entry of the table is $q=566999$.

\section{The search}\label{sec:algo}

\subsection{Problem and notation}
Fix an odd constant $c$ and a residue $r\in\{1,5\}$. Let
$\mathcal S=\{2^x+c\,2^y: x,y\ge 1\}$ and let $\mathcal P_r$ be the set of primes $\equiv r\pmod 6$.
We must list all odd $m<N$ with $m\notin \mathcal P_r+\mathcal S$. The set $\mathcal S$ is very thin:
it has about $(\log_2 N)^2$ elements below $N$ (about $1500$ for $N=10^{12}$), while a candidate
$m-s$ is prime with probability of order $1/\log N$. A direct test of each $m$ therefore needs a dozen
or more primality tests on average and, more importantly, hundreds for the rare $m$ that matter.

\subsection{Algorithm}
We store the indicator of $\mathcal P_r$ on odd numbers as a bitmap $P$, bit $g$ standing for $2g+1$.
For odd $m=2j+1$ and even $s\in\mathcal S$ the number $m-s$ corresponds to bit $j-s/2$. Hence the
bitmap $C$ of represented numbers is an OR of shifted copies of $P$:
\[
   C \;=\; \bigvee_{s\in\mathcal S} \bigl(P \gg s/2\bigr).
\]
The program proceeds in segments of $2^{21}$ bits ($256$\,KB, so that a segment and the data it
reads stay in the second-level cache):
\begin{enumerate}
\item \emph{Sieve.} A segmented sieve of Eratosthenes on odd numbers produces $P$ on the segment;
a periodic mask keeps only the residue class $r$ modulo $6$. A window of the previous $2^{26}$ bits of
$P$ is kept.
\item \emph{Bit-parallel cover.} For the $K=256$ smallest elements $s\in\mathcal S$ the shifted
bitmap is OR-ed into $C$, $64$ numbers per machine word.
\item \emph{Residue.} Each $m$ whose bit in $C$ is still zero is treated individually: the
remaining $s\in\mathcal S$ are tried in increasing order, by a single bit lookup while $s/2$ lies inside the
window and by a deterministic Miller--Rabin test (twelve prime bases, valid below $2^{64}$) otherwise.
An $m$ for which no $s$ succeeds is reported.
\end{enumerate}
Steps 1--3 decide membership in $\mathcal P_r+\mathcal S$ exactly; no probabilistic step is involved.

\subsection{Cost}
Let $M=N/2$ be the number of odd integers tested. The sieve costs $O(M\log\log N)$ bit operations
and uses the $O(\sqrt N/\log N)$ primes below $\sqrt N$. The cover costs $KM/64$ word operations.
If $u_K$ denotes the proportion of $m$ not covered by the first $K$ shifts, the residue step costs
at most $u_K M\,|\mathcal S\cap[1,N]|$ lookups. Each shift covers a given $m$ with probability about
$\delta\approx 3/\log N$ when $m-s$ lies in the right class modulo~$6$ (one shift in four or one in
two, depending on $m \bmod 6$), so $u_K$ decays geometrically in $K$; for $N=10^{12}$ and $K=256$
we measured $u_K\approx 7\cdot 10^{-4}$, and about five Miller--Rabin tests per million integers.
Memory use is $16$\,MB per process. Independent ranges are run in parallel.

On a six-core desktop processor (AMD Ryzen 5 7600X) the program tests about $4\cdot 10^{9}$ integers
per second; each part of Theorem~\ref{thm:search} took four to six minutes.

\subsection{Validation}
The statement of Theorem~\ref{thm:search} depends on the correctness of the program, so we record
the checks that were made.
\begin{itemize}
\item For each sequence, the list of unrepresentable $m<2\cdot 10^{5}$ produced by the program equals
the list produced by a brute-force program written directly from the definition, also when the
program is run with a tiny window and $K=3$, which forces almost all work through step~3.
\item The brute-force representation counts reproduce the terms published in the OEIS entries
A157218, A157225, A157237 and A157242.
\item The program finds Adamczewski's counterexample $1433987797$ for A157225, and for a related
conjecture of Sun and Hou on Pell numbers (A154364) it finds the four counterexamples near
$1.7\cdot 10^{11}$ published by McNeil in 2009.
\item The run for part (i) was produced by an earlier version of the program with a different
memory layout; both versions agree on all ranges where both were run.
\item Each reported counterexample was confirmed by three programs sharing no code with the search:
a Python script using its own Miller--Rabin test, a Python script using \texttt{sympy}, and PARI/GP's
\texttt{isprime}.
\end{itemize}
The negative statements (``no other counterexample'') have not been reproduced by an independent
implementation; we invite such a check.

\section{Why this number?}\label{sec:structure}

\subsection{Local densities}
The residue of $m$ modulo $6$ decides which parities of $(x,y)$ can occur. For A157218:
\begin{center}
\begin{tabular}{lll}
\toprule
$m \bmod 6$ & parities of $(x,y)$ & candidates \\
\midrule
$5$ & both odd & $m-2\cdot 4^a-14\cdot 4^b$, \ $a,b\ge 0$ \\
$3$ & both even & $m-4^a-7\cdot 4^b$, \ $a,b\ge 1$ \\
$1$ & opposite & twice as many candidates \\
\bottomrule
\end{tabular}
\end{center}
(For A157225 the three rows are $m\equiv 3,1,5 \pmod 6$ respectively.) Counterexamples should
therefore come from the first two classes, and all five known ones do.

Fix one of the first two classes, with candidates $m-A\,4^a-B\,4^b$, and a prime $q\ge 5$. Let
$d=\ord_q(4)$. Whether $q$ divides a candidate depends only on $(a \bmod d,\, b\bmod d)$, so the
proportion of candidates divisible by $q$ is
\[
  \rho_q(m) \;=\; \frac{1}{d^{2}}\,\#\bigl\{(a,b)\in(\mathbb Z/d)^2 : A\,4^a+B\,4^b\equiv m \pmod q\bigr\},
\]
a function of $m \bmod q$. For a random integer this proportion would be $1/q$. The prime $q$ thus
multiplies the expected number of prime candidates by $f_q(m)=(1-\rho_q(m))/(1-1/q)$, and we set
\[
   S(m)=\prod_{q\ge 5} f_q(m),\qquad T(m)=\sum_{(x,y)}\frac{3}{\log(m-2^x-7\cdot 2^y)},
\]
the sum being over the admissible pairs; $3/\log n$ is the density of primes among integers
$\equiv 1\pmod 6$ near $n$. Heuristically $m$ has about $E(m)=S(m)\,T(m)$ representations.

The factor $\rho_q$ can be much larger than $1/q$ when $4$ has small order modulo $q$. If $q\mid m$
and $-B/A$ is a power of $4$ modulo $q$, then $q$ divides the candidate exactly when $b-a$ lies in
one residue class modulo $d$, so $\rho_q=1/d$. This is what happens for $\mz$ and $q=11$:
$11\mid \mz$, $d=5$, and one candidate in five is divisible by $11$ instead of one in eleven.

Table~\ref{tab:local} shows that $\mz$ lies in the least favourable residue class for $5$, $7$,
$11$, $13$, $31$ and $41$ simultaneously. The product over all primes used in our computation is
$S(\mz)\approx 0.615$, a value that only about $1\%$ of the integers $m\equiv 5\pmod 6$ attain
(the median is $1.02$; the smallest value over all residue classes is about $0.29$). With
$T(\mz)=38.4$ this gives $E(\mz)\approx 23.6$ expected representations in place of $38$. After
removing candidates with a prime factor below $500$ only $41$ of the $342$ remain, and all of them
turn out to be composite.

\begin{table}[t]
\centering
\begin{tabular}{rrrrrr}
\toprule
$q$ & $\ord_q(4)$ & $\mz \bmod q$ & $\rho_q(\mz)$ & $\max_m \rho_q(m)$ & $1/q$ \\
\midrule
5 & 2 & 1 & 0.2500 & 0.2500 & 0.2000 \\
7 & 3 & 1 & 0.3333 & 0.3333 & 0.1429 \\
11 & 5 & 0 & 0.2000 & 0.2000 & 0.0909 \\
13 & 6 & 3 & 0.0833 & 0.0833 & 0.0769 \\
17 & 4 & 11 & 0.0625 & 0.1250 & 0.0588 \\
31 & 5 & 15 & 0.0800 & 0.0800 & 0.0323 \\
41 & 10 & 27 & 0.0400 & 0.0400 & 0.0244 \\
\bottomrule
\end{tabular}
\caption{Proportion $\rho_q$ of the candidates $\mz-2\cdot 4^a-14\cdot 4^b$ divisible by $q$,
compared with the largest value over all residue classes of $m$ and with the proportion $1/q$ for
random integers.}\label{tab:local}
\end{table}

\subsection{How many counterexamples?}
Treating the candidates as independent, the probability that $m$ has no representation is about
\begin{equation}\label{eq:model}
  \prod_{(x,y)}\Bigl(1-\frac{3\,S(m)}{\log(m-2^x-7\cdot 2^y)}\Bigr).
\end{equation}
Summing \eqref{eq:model} over $m$, with $m \bmod q$ uniform and independent for different $q$,
gives the expected number of counterexamples in a range. Table~\ref{tab:pred} compares this with
the search. The model is the same for A157218 and A157225, which differ only in the assignment of
the two parity classes to residues of $m$ modulo~$6$.

\begin{table}[t]
\centering
\begin{tabular}{lccc}
\toprule
 & \multicolumn{2}{c}{$1.5\cdot 10^{8}\le m<10^{12}$} & $10^{12}\le m<10^{13}$\\
class of $(x,y)$ & predicted & found (A157218 / A157225) & predicted \\
\midrule
both odd & 0.34 & 1 \ / \ 1 & 0.13 \\
both even & 1.27 & 0 \ / \ 3 & 0.46 \\
\midrule
total per conjecture & 1.6 & 1 \ / \ 4 & 0.59 \\
\bottomrule
\end{tabular}
\caption{Expected number of counterexamples under the model \eqref{eq:model}, per conjecture,
and the numbers found. The five counterexamples are $445466268731$ (both odd) for A157218, and
$367084269651$ (both odd), $1433987797$, $45483661477$, $85305551767$ (both even) for A157225.}
\label{tab:pred}
\end{table}

\begin{heuristic}\label{heur}
Under the model \eqref{eq:model}, the expected number of counterexamples to A157218 in
$[10^{k},10^{k+1})$ is $0.6$ for $k=12$, $0.7$ for $k=13$, about one per decade up to $10^{16}$,
and it keeps growing. In particular the model predicts infinitely many counterexamples.
\end{heuristic}

Two warnings are in order. First, \eqref{eq:model} is exponentially sensitive to second-order
choices: replacing the product by $\exp(-S(m)T(m))$ multiplies the predictions by about four, while
conditioning on the absence of prime factors below $100$ divides them by more than twenty. The
agreement in Table~\ref{tab:pred} should be read as order-of-magnitude support, no more. Second,
the model says nothing rigorous.

\begin{question}
Are there infinitely many odd integers not of the form $p+2^x+7\cdot 2^y$ with $p\equiv 1 \pmod 6$
prime? Is there one in every interval $[X,2X]$ for large $X$?
\end{question}

Crocker's proof for $p+2^a+2^b$ uses the factorisation of $2^a+2^b$ and the divisibility of
$2^{2^n}-1$ by Fermat numbers; we see no analogue for $2^x+7\cdot 2^y$. A covering-congruence proof
in the style of Erd\H{o}s would require finitely many primes $q$ and residues $m \bmod q$ whose
exceptional sets cover all pairs $(a,b)$; since each prime covers a proportion of at most about
$1/\ord_q(4)$ of the pairs, and the pairs covered by different primes overlap, we do not know
whether such a covering exists.

\section*{Use of AI and acknowledgements}
The search program, the analysis scripts and the text of this note were written with the help of
an AI assistant (Claude, Anthropic), which also proposed the family of conjectures to test. All
statements are backed by programs in the repository cited above, and the counterexample itself can
be checked independently of them from Table~\ref{tab:cert}. We thank Zhi-Wei Sun for the
conjectures.

\begin{thebibliography}{9}
\bibitem{crocker} R.~Crocker, \emph{On the sum of a prime and of two powers of two}, Pacific J. Math.
\textbf{36} (1971), 103--107.
\bibitem{erdos} P.~Erd\H{o}s, \emph{On integers of the form $2^k+p$ and some related problems},
Summa Brasil. Math. \textbf{2} (1950), 113--123.
\bibitem{oeisopen} T.~Adamczewski, \emph{OEIS Open: How many conjectures can language models turn into
theorems?}, arXiv:2608.11941 (2026).
\bibitem{oeis} OEIS Foundation Inc., \emph{The On-Line Encyclopedia of Integer Sequences},
\url{https://oeis.org}, sequences A157218, A157225, A157237, A154364.
\bibitem{pan} H.~Pan, \emph{On the integers not of the form $p+2^a+2^b$}, Acta Arith. \textbf{148}
(2011), 55--61.
\bibitem{sun} Z.-W.~Sun, \emph{Mixed sums of primes and other terms}, arXiv:0901.3075 (2009).
\end{thebibliography}

\appendix
\section{Certificate for Theorem~\ref{thm:main}}

For each pair of odd $x,y$ with $2^x+7\cdot 2^y<\mz$, Table~\ref{tab:cert} gives the least prime
factor $q$ of $\mz-2^x-7\cdot 2^y$. The table is produced by \texttt{verify.py --certificate} in the
repository and can be checked with $342$ divisions.

\begin{table}[p]
\centering
\scriptsize
\setlength{\tabcolsep}{3pt}
% generated by make_tables.py -- do not edit
\begin{tabular}{rrr@{\qquad}rrr@{\qquad}rrr@{\qquad}rrr@{\qquad}rrr@{\qquad}rrr}
\toprule
$x$ & $y$ & $q$ & $x$ & $y$ & $q$ & $x$ & $y$ & $q$ & $x$ & $y$ & $q$ & $x$ & $y$ & $q$ & $x$ & $y$ & $q$ \\
\midrule
1 & 1 & 5 & 7 & 7 & 29 & 13 & 13 & 5 & 19 & 19 & 103 & 25 & 25 & 5 & 31 & 31 & 171679 \\
1 & 3 & 496187 & 7 & 9 & 73 & 13 & 15 & 31 & 19 & 21 & 17 & 25 & 27 & 73 & 31 & 33 & 1831 \\
1 & 5 & 5 & 7 & 11 & 97 & 13 & 17 & 5 & 19 & 23 & 29 & 25 & 29 & 5 & 31 & 35 & 23 \\
1 & 7 & 3061 & 7 & 13 & 227 & 13 & 19 & 229 & 19 & 25 & 89 & 25 & 31 & 29 & 33 & 1 & 5 \\
1 & 9 & 5 & 7 & 15 & 11 & 13 & 21 & 5 & 19 & 27 & 11 & 25 & 33 & 5 & 33 & 3 & 7 \\
1 & 11 & 141403 & 7 & 17 & 727 & 13 & 23 & 251 & 19 & 29 & 17 & 25 & 35 & 197 & 33 & 5 & 5 \\
1 & 13 & 5 & 7 & 19 & 23 & 13 & 25 & 5 & 19 & 31 & 157 & 27 & 1 & 7 & 33 & 7 & 7 \\
1 & 15 & 461 & 7 & 21 & 3167 & 13 & 27 & 41 & 19 & 33 & 53 & 27 & 3 & 7 & 33 & 9 & 5 \\
1 & 17 & 5 & 7 & 23 & 109 & 13 & 29 & 5 & 19 & 35 & 661 & 27 & 5 & 7 & 33 & 11 & 7 \\
1 & 19 & 11 & 7 & 25 & 11 & 13 & 31 & 11 & 21 & 1 & 5 & 27 & 7 & 7 & 33 & 13 & 5 \\
1 & 21 & 5 & 7 & 27 & 73 & 13 & 33 & 5 & 21 & 3 & 7 & 27 & 9 & 7 & 33 & 15 & 7 \\
1 & 23 & 113 & 7 & 29 & 89 & 13 & 35 & 31 & 21 & 5 & 5 & 27 & 11 & 7 & 33 & 17 & 5 \\
1 & 25 & 5 & 7 & 31 & 37 & 15 & 1 & 7 & 21 & 7 & 7 & 27 & 13 & 7 & 33 & 19 & 7 \\
1 & 27 & 2711 & 7 & 33 & 136501 & 15 & 3 & 7 & 21 & 9 & 5 & 27 & 15 & 7 & 33 & 21 & 5 \\
1 & 29 & 5 & 7 & 35 & 11 & 15 & 5 & 7 & 21 & 11 & 7 & 27 & 17 & 7 & 33 & 23 & 7 \\
1 & 31 & 79 & 9 & 1 & 5 & 15 & 7 & 7 & 21 & 13 & 5 & 27 & 19 & 7 & 33 & 25 & 5 \\
1 & 33 & 5 & 9 & 3 & 7 & 15 & 9 & 7 & 21 & 15 & 7 & 27 & 21 & 7 & 33 & 27 & 7 \\
1 & 35 & 2143 & 9 & 5 & 5 & 15 & 11 & 7 & 21 & 17 & 5 & 27 & 23 & 7 & 33 & 29 & 5 \\
3 & 1 & 7 & 9 & 7 & 7 & 15 & 13 & 7 & 21 & 19 & 7 & 27 & 25 & 7 & 33 & 31 & 7 \\
3 & 3 & 7 & 9 & 9 & 5 & 15 & 15 & 7 & 21 & 21 & 5 & 27 & 27 & 7 & 33 & 33 & 5 \\
3 & 5 & 7 & 9 & 11 & 7 & 15 & 17 & 7 & 21 & 23 & 7 & 27 & 29 & 7 & 33 & 35 & 7 \\
3 & 7 & 7 & 9 & 13 & 5 & 15 & 19 & 7 & 21 & 25 & 5 & 27 & 31 & 7 & 35 & 1 & 31 \\
3 & 9 & 7 & 9 & 15 & 7 & 15 & 21 & 7 & 21 & 27 & 7 & 27 & 33 & 7 & 35 & 3 & 11 \\
3 & 11 & 7 & 9 & 17 & 5 & 15 & 23 & 7 & 21 & 29 & 5 & 27 & 35 & 7 & 35 & 5 & 17 \\
3 & 13 & 7 & 9 & 19 & 7 & 15 & 25 & 7 & 21 & 31 & 7 & 29 & 1 & 5 & 35 & 7 & 29 \\
3 & 15 & 7 & 9 & 21 & 5 & 15 & 27 & 7 & 21 & 33 & 5 & 29 & 3 & 360907 & 35 & 9 & 13 \\
3 & 17 & 7 & 9 & 23 & 7 & 15 & 29 & 7 & 21 & 35 & 7 & 29 & 5 & 5 & 35 & 11 & 23 \\
3 & 19 & 7 & 9 & 25 & 5 & 15 & 31 & 7 & 23 & 1 & 11 & 29 & 7 & 11 & 35 & 13 & 11 \\
3 & 21 & 7 & 9 & 27 & 7 & 15 & 33 & 7 & 23 & 3 & 421 & 29 & 9 & 5 & 35 & 15 & 87037 \\
3 & 23 & 7 & 9 & 29 & 5 & 15 & 35 & 7 & 23 & 5 & 31 & 29 & 11 & 13 & 35 & 17 & 83 \\
3 & 25 & 7 & 9 & 31 & 7 & 17 & 1 & 5 & 23 & 7 & 71 & 29 & 13 & 5 & 35 & 19 & 2267 \\
3 & 27 & 7 & 9 & 33 & 5 & 17 & 3 & 297833 & 23 & 9 & 13 & 29 & 15 & 311 & 35 & 21 & 13 \\
3 & 29 & 7 & 9 & 35 & 7 & 17 & 5 & 5 & 23 & 11 & 11 & 29 & 17 & 5 & 35 & 23 & 11 \\
3 & 31 & 7 & 11 & 1 & 271517 & 17 & 7 & 476849 & 23 & 13 & 397 & 29 & 19 & 23 & 35 & 25 & 3643 \\
3 & 33 & 7 & 11 & 3 & 23 & 17 & 9 & 5 & 23 & 15 & 19 & 29 & 21 & 5 & 35 & 27 & 191 \\
3 & 35 & 7 & 11 & 5 & 17 & 17 & 11 & 13 & 23 & 17 & 229 & 29 & 23 & 13 & 35 & 29 & 17 \\
5 & 1 & 5 & 11 & 7 & 101 & 17 & 13 & 5 & 23 & 19 & 32801 & 29 & 25 & 5 & 35 & 31 & 31 \\
5 & 3 & 11 & 11 & 9 & 11 & 17 & 15 & 11 & 23 & 21 & 11 & 29 & 27 & 11 & 35 & 33 & 11 \\
5 & 5 & 5 & 11 & 11 & 183487 & 17 & 17 & 5 & 23 & 23 & 199 & 29 & 29 & 5 & 35 & 35 & 29 \\
5 & 7 & 2027 & 11 & 13 & 17 & 17 & 19 & 223 & 23 & 25 & 31 & 29 & 31 & 109 & 37 & 1 & 5 \\
5 & 9 & 5 & 11 & 15 & 41983 & 17 & 21 & 5 & 23 & 27 & 24439 & 29 & 33 & 5 & 37 & 3 & 88969 \\
5 & 11 & 13 & 11 & 17 & 45161 & 17 & 23 & 13 & 23 & 29 & 269 & 29 & 35 & 13 & 37 & 5 & 5 \\
5 & 13 & 5 & 11 & 19 & 11 & 17 & 25 & 5 & 23 & 31 & 11 & 31 & 1 & 2417 & 37 & 7 & 532489 \\
5 & 15 & 19 & 11 & 21 & 13 & 17 & 27 & 37 & 23 & 33 & 13 & 31 & 3 & 19 & 37 & 9 & 5 \\
5 & 17 & 5 & 11 & 23 & 773 & 17 & 29 & 5 & 23 & 35 & 31 & 31 & 5 & 283 & 37 & 11 & 4943 \\
5 & 19 & 59 & 11 & 25 & 23 & 17 & 31 & 71 & 25 & 1 & 5 & 31 & 7 & 80147 & 37 & 13 & 5 \\
5 & 21 & 5 & 11 & 27 & 184231 & 17 & 33 & 5 & 25 & 3 & 11 & 31 & 9 & 11 & 37 & 15 & 11 \\
5 & 23 & 11 & 11 & 29 & 11 & 17 & 35 & 11 & 25 & 5 & 5 & 31 & 11 & 83 & 37 & 17 & 5 \\
5 & 25 & 5 & 11 & 31 & 349 & 19 & 1 & 331 & 25 & 7 & 566999 & 31 & 13 & 23 & 37 & 19 & 1061 \\
5 & 27 & 38791 & 11 & 33 & 13 & 19 & 3 & 89 & 25 & 9 & 5 & 31 & 15 & 90001 & 37 & 21 & 5 \\
5 & 29 & 5 & 11 & 35 & 61 & 19 & 5 & 17 & 25 & 11 & 31 & 31 & 17 & 809 & 37 & 23 & 1051 \\
5 & 31 & 31 & 13 & 1 & 5 & 19 & 7 & 11 & 25 & 13 & 5 & 31 & 19 & 11 & 37 & 25 & 5 \\
5 & 33 & 5 & 13 & 3 & 19 & 19 & 9 & 157477 & 25 & 15 & 101 & 31 & 21 & 19 & 37 & 27 & 149 \\
5 & 35 & 13 & 13 & 5 & 5 & 19 & 11 & 61 & 25 & 17 & 5 & 31 & 23 & 131 & 37 & 29 & 5 \\
7 & 1 & 151 & 13 & 7 & 37 & 19 & 13 & 17 & 25 & 19 & 2113 & 31 & 25 & 21019 & 37 & 31 & 151 \\
7 & 3 & 599 & 13 & 9 & 5 & 19 & 15 & 2753 & 25 & 21 & 5 & 31 & 27 & 271 & 37 & 33 & 5 \\
7 & 5 & 11 & 13 & 11 & 11 & 19 & 17 & 11 & 25 & 23 & 11 & 31 & 29 & 11 & 37 & 35 & 11 \\
\bottomrule
\end{tabular}

\caption{Least prime factor $q$ of $445466268731-2^x-7\cdot 2^y$ for all $342$ pairs of odd $x,y$.}
\label{tab:cert}
\end{table}

\end{document}
